\documentclass[11pt,twoside,openright]{memoir}
% Idioma predeterminado de toda la colección: español. Cambiar a english si hace falta.
\def\AviaLanguage{spanish}
% Perfil narrativo: columna amplia y notas a pie. Perfil ilustrado: márgenes laterales.
\newif\ifAviaIllustrated
\AviaIllustratedfalse % perfil narrativo comercial por defecto
% A 6x9 pulgadas: 152.4 x 228.6 mm. Perfil ilustrado sacrifica ancho de texto para margen.
\setstocksize{9in}{6in}
\settrimmedsize{9in}{6in}{*}
\settrims{0pt}{0pt}
\ifAviaIllustrated
  \settypeblocksize{7.2in}{3.55in}{*}
  \setlrmargins{0.67in}{*}{*}
  \setmarginnotes{10pt}{0.9in}{10pt}
\else
  \settypeblocksize{7.15in}{4.35in}{*}
  \setlrmargins{0.88in}{*}{*}
  \setmarginnotes{8pt}{0.3in}{10pt}
\fi
\setulmargins{0.85in}{*}{*}
\checkandfixthelayout
% Impresión a doble cara; verificar margen de encuadernación según número de páginas e imprenta.

\usepackage{avia-editorial}
\usepackage{amsmath,amssymb,longtable,booktabs,array,tikz}
\usetikzlibrary{arrows.meta,positioning,babel}
\usepackage[backend=biber,style=authoryear,sorting=nyt]{biblatex}
\addbibresource{references.bib}
\AtBeginBibliography{\raggedright}
\title{Avia Editorial}
\booksubtitle{Edición 2.2\\Lectura, ilustraciones y pruebas de composición}
\author{Una plantilla para escribir en español}
\hypersetup{pdftitle={Avia Editorial 2.2},pdfauthor={Avia Editorial}}
\begin{document}
\frontmatter
\aviatitlepage
\chapter*{Antes de empezar}
Esta muestra es un documento editorial y técnico, no un manual de aviación. Enseña los recursos del diseño con contenido original y separa la demostración de las pruebas sintéticas. Los archivos de pruebas fuerzan situaciones límite; no deben confundirse con capítulos de un libro ni contarse como texto pedagógico.

El perfil narrativo es el predeterminado. Las notas aparecen al pie y la columna principal aprovecha el ancho de lectura. El perfil ilustrado reserva espacio para notas breves. Ambos mantienen el interior monocromático, la apertura de capítulos y la paginación exterior.

\begin{insightbox}[Criterio de aceptación]
Un proceso de compilación que termina no demuestra por sí solo que el PDF sea correcto. La entrega pasa también por controles de fuentes, referencias, contenido recortado, color y casos de regresión. Cada imprenta y cada libro necesitan, además, una comprobación editorial y física propias.
\end{insightbox}
\tableofcontents*
\mainmatter
\chapter{La curiosidad empieza en tierra}
\chapterquote{Aprender comienza cuando una pregunta reemplaza una certeza.}{Nota editorial}
\opening{Una cabina pequeña puede abrir preguntas enormes.}

Este es un capítulo de demostración, no un procedimiento de vuelo. Un libro de divulgación puede alternar relatos personales, observaciones y explicaciones técnicas sin convertirse en un manual operacional.

Aquí aparece una nota al margen\sidenote{Las notas marginales son ideales para una definición corta, un dato o una referencia que no interrumpa el relato.}. La composición está inspirada en la tradición editorial de Edward Tufte, pero construida sobre \texttt{memoir}.

\section{La diferencia entre mirar y comprender}

La pregunta interesante no es solo qué hace una pieza, sino qué problema resuelve. Esta clase de libro se beneficia de una lectura continua y de ilustraciones puntuales.

\begin{insightbox}[Idea para recordar]
Una buena nota visual debe aclarar el contenido principal, no competir con él.
\end{insightbox}

\margintext{Una idea lateral puede leerse por separado, como segunda capa de contenido.}Un texto más largo permite comprobar el uso de notas y el espaciado entre líneas. La plantilla acepta capítulos independientes y bibliografía administrada por BibLaTeX. Una referencia de ejemplo puede citarse con \textcite{tufte2001}.

\section{Una imagen que explique}

\aviamarginfigure{assets/wing.pdf}{Esquema simplificado, dibujado para demostrar una figura marginal.}La figura lateral no necesita ocupar una página completa. Puede acompañar una explicación breve sin interrumpir el párrafo.

\begin{fullwidth}
\begin{insightbox}[Lectura extendida]
Este recuadro ocupa el ancho del texto y el margen lateral. Úsalo para conceptos relevantes, ilustraciones, resúmenes o fotografías panorámicas. En libros de impresión a doble cara es importante probar el diseño de márgenes interiores y exteriores.
\end{insightbox}
\end{fullwidth}

\chapter{Pensar en sistemas}
\opening{Comprender un avión es comprender relaciones entre componentes.}

En un capítulo técnico conviene separar narración y recuadros didácticos.

\begin{technicalbox}[Precaución editorial]
Esta obra es una guía educativa no oficial. No sustituye el manual de vuelo aprobado, las listas de verificación aplicables ni la instrucción aeronáutica certificada.
\end{technicalbox}

\section{Un ejemplo de ilustración}

\widefigure[0.8\linewidth]{assets/wing.pdf}{Esquema conceptual de demostración, no apto para instrucción operacional.}

La ilustración se guarda como archivo independiente y se reutiliza en distintos capítulos. La herramienta se puede adaptar a libros narrativos con ilustraciones ocasionales, manuales ilustrados y guías comerciales.

\sidecite[Una nota bibliográfica editorial]{Tufte, 2001}Una referencia lateral ayuda a ofrecer contexto sin convertir cada página en un artículo académico.

\section{Una idea final}

La plantilla distingue entre el contenido, que puede crecer cientos de páginas, y las decisiones de maquetación, que pueden cambiar antes de imprimir.

\chapter{De la muestra a tu libro}
\section{Escritura y compilación}
El archivo principal define el título, el autor y el orden de capítulos. La configuración selecciona idioma y perfil. El estilo editorial contiene las macros reutilizables. Esta separación evita copiar correcciones de un capítulo a otro y permite cambiar el diseño sin reescribir el manuscrito.

El flujo recomendado crea una carpeta de compilación limpia, conserva los registros y verifica el PDF antes de copiarlo a la carpeta de entregas aprobadas. Un fallo de la compilación actual no debe dejar disponible un PDF antiguo que parezca nuevo. Por eso el directorio de aprobación se vacía al comenzar cada ejecución.

\section{Qué prueban las simulaciones}
Las pruebas positivas incluyen recuadros extensos, ilustraciones altas, pies largos, notas al pie, notas marginales, páginas enfrentadas, tablas de varias páginas y un documento sintético de cientos de páginas. Los marcadores permiten comprobar que el texto inicial y final de cada caso no desaparece al componerlo.

Las pruebas negativas introducen defectos de forma deliberada. Una referencia inexistente, un carácter sin glifo, un objeto fuera de la página o una nota marginal excesiva deben impedir la aprobación. Rechazar esos ejemplos es el resultado correcto: no se incluyen entre los PDF aprobados.

\begin{technicalbox}[Un límite no es un fallo oculto]
Los entornos de ancho extendido, las imágenes y las notas marginales son objetos con restricciones físicas. La versión 2.2 comprueba sus dimensiones y detiene usos incompatibles. Las cajas normales, en cambio, admiten contenido divisible. No se ocultan las advertencias del motor para obtener una marca verde.
\end{technicalbox}

\section{La revisión que no se automatiza}
Un control geométrico no demuestra exactitud de una ilustración, calidad de un argumento ni legibilidad de una fotografía en un papel concreto. La prueba física y la revisión de contenido siguen siendo necesarias. El informe registra qué se ejecutó y con qué versiones; no certifica todas las combinaciones de paquetes que puedan añadirse después.

Al crear un nuevo libro, conserva las pruebas del proyecto. Su función no es decorar el repositorio: permiten detectar cuando un cambio vuelve a introducir una pérdida de contenido o una extensión hacia el lado incorrecto de la hoja. La tabla~\ref{tab:control} resume qué responsabilidad corresponde a cada revisión.

\backmatter
\printbibliography[heading=bibintoc,title={Referencias del ejemplo}]
\end{document}
